\documentclass{ctexart}

\usepackage{graphicx}
\usepackage{setspace}
\usepackage{booktabs}
\usepackage{amsmath} 
\usepackage{tcolorbox}

\title{考研英语笔记}
\author{杨子颉}

\begin{document}
\maketitle
\newpage
\tableofcontents
\newpage

\section{名词词组}
   \subsection{名词词组的组成}
   名词出来必有标记
    \begin{tcolorbox}[colback=blue!5!white]
        名词出来必有标记

       名词=限定词+修饰词+名词

       名词：名词有时候可以省略，重复的时候名词可以省略

       限定词：要么和指代有关、要么和数量有关

       修饰词：一般由形容词充当，数量不做要求，但限定词一定要有
    \end{tcolorbox}
    \subsection{名词的分类}
        \subsubsection{普通名词}
        \begin{itemize}
            \item[1)] 普通名词：根据语境决定单复数，比如family表示家庭的时候用单数，表示家人的时候用复数
            \item[2)] 专有名词：具有唯一性和特指性，但是可以和普通名词相互转换，比如班长，一个班只有一个班长，这时候是专有名词 
        \end{itemize}
        \subsubsection{限定词}
        对名词词组中的中心词起\textbf{特指、类指、确定或非确定数量}等限定作用（专有名词不需要限定）

        比如the、an/a（还有零冠词）、你的我的他的、一半、全部.....等
            \begin{itemize}
                \item[1)] 前位限定词：表全部、倍数、分数
                \item[2)] 中位限定词：表所属、范围限定（你的、我的、他的...）（大部分为冠词）
                \item[3)] 后位限定词：表多少、表数字
                \item[4)] 当多个限定词修饰一个名词的时候，顺序为\textbf{前＞中＞后}但并非都包含三个
                \item[5)] 零冠词：
                    \begin{itemize}
                        \item[(1)] 指名词前不加冠词和定冠词，主要用于\textbf{泛指通用}
                        \item[1)] 可数名词的复数与不可数名词表示泛指通用的时候可以零冠词
                    \end{itemize}
            \end{itemize}
        \subsubsection{修饰词}
        名词词组中由形容词充当，数量并没有要求，没有或多个都可以

        名词词组做修饰词要用连字符



\end{document}